\section{Work Importance}

This thesis addresses fundamental limitations in current ETA prediction research by developing a comprehensive evaluation framework and validating route-aware dynamic graph neural networks on real-world data.

\textbf{Research Community Impact:} Current ETA prediction research suffers from fragmentation, where different model architectures are evaluated on incompatible datasets, making fair comparison impossible. This thesis addresses this gap by developing a real-world dataset enhancement approach that enables rigorous evaluation of route-aware dynamic graph neural networks on actual traffic conditions while preserving ground truth travel times, providing the research community with a validated approach for evaluating advanced ETA prediction models.

\textbf{Theoretical and Methodological Contribution:} Initial results from the DSTRA-GNN model (detailed in Appendix A) on simulated data demonstrate that explicit route information combined with vehicle-level dynamic graph modeling significantly improves ETA prediction accuracy (46.2\,s MAE, 82.3\% improvement over average baseline). This thesis will provide the first comprehensive validation of route-aware dynamic graph neural networks on enhanced real-world datasets, establishing whether theoretical advantages translate to practical improvements. The real-world dataset enhancement approach enables evaluation of advanced models on actual traffic conditions while preserving ground truth travel times, establishing a new paradigm combining real data realism with structural richness required by modern architectures.

\textbf{Practical Value:} Accurate ETA prediction directly impacts navigation systems, ride-sharing platforms, fleet management, and logistics operations. By validating route-aware dynamic graph models on real-world data, this research provides evidence-based guidance for deploying advanced ETA prediction systems in production environments, with potential benefits for millions of daily users of transportation services.

